% UBT-AI-PROVENANCE-BEGIN
% schema: ubt-ai-provenance/v1
% tier: C_working
% ai_assistance: disclosed
% human_review: risk-based
% editorial_responsibility: Ing. David Jaroš
% policy: ../../../AI_PROVENANCE.md
% notice: Working material; exhaustive human review is not claimed.
% UBT-AI-PROVENANCE-END

\chapter{The fine structure constant program: what is derived and what is not}
\label{ch:alpha}

\section{Status before algebra}
The authoritative repository status is deliberately stronger than any elegant
formula in this chapter:
\begin{quote}
The physical fine-structure constant is not yet derived from first principles
in UBT.  The integer-$137$ route is structural/conditional, while the bridge
fixing the required effective coefficient from the UBT action remains open.
\end{quote}
The live sources of status are \texttt{STATUS.md} and
\texttt{WHAT\_IS\_PROVED.md}. The alpha ledger is\\
\texttt{canonical/alpha/}\\
\texttt{ALPHA\_MASTER\_STATUS.md}. Historical texts using the
phrase ``fit-free alpha derivation'' do not override these ledgers.

\section{Conditional winding potential}
A central reduced model studies
\[
 \boxed{V_{\rm eff}(n)=n^2-Bn\ln n,\qquad n>0.}
\]
Treating $n$ temporarily as a continuous variable,
\[
 \frac{dV_{\rm eff}}{dn}
 =2n-B(\ln n+1).
\]
A stationary point therefore satisfies
\[
 \boxed{B=\frac{2n}{\ln n+1}.}
\]
If one asks for a stationary point at $n=137$, the required coefficient is
\[
 B_{137}=\frac{274}{\ln137+1}\approx46.284.
\]
This equation is an exact property of the reduced potential.  It does
\emph{not} explain why the UBT action must generate that value of $B$.

The second derivative is
\[
 \frac{d^2V_{\rm eff}}{dn^2}=2-\frac{B}{n}.
\]
At a stationary point this becomes
\[
 2-\frac{2}{\ln n+1}
 =\frac{2\ln n}{\ln n+1}>0
 \qquad(n>1),
\]
so the continuous stationary point is a local minimum.  Integer and
prime-sector restrictions can then be studied as an additional discrete
selection problem.

\section{What ``prime stability'' can and cannot mean}
The repository contains a structural prime-stability analysis in which the
integer sector around $137$ is examined under a specified selection rule.  A
prime-labelled minimum can be a reproducible result of that reduced model.
However:
\begin{itemize}
 \item the primality of $137$ is standard number theory;
 \item selection of a prime sector does not by itself identify
       $\alpha^{-1}$ with that sector;
 \item the physical correction from $137$ to $137.036\ldots$ needs a separate
       derivation;
 \item the coefficient $B$ must still be obtained from the UBT dynamics rather
       than inferred from the target value.
\end{itemize}
The theta/Mellin/Feynman chapter develops possible mechanisms for understanding
where multiplicative prime structure could enter, but those mechanisms are
research bridges, not completed alpha proofs.

\section{Audit \(N_{\rm eff}\): two different counts must not be confused}
A previous derivation used
\[
 N_{\rm eff}=3\times2\times2=12.
\]
The current audit distinguishes at least two objects:
\begin{align*}
 N_{\rm eff}^{\rm twist}&=12,
 &&(\mathrm T),\\
 N_{\rm eff}^{\rm loop}&=3,
 &&(\mathrm L).
\end{align*}
Here $\mathrm T$ denotes the SU(2)-twist construction and $\mathrm L$ the
direct audited scalar-action loop count.
The identification
\[
 N_{\rm eff}^{\rm loop}=N_{\rm eff}^{\rm twist}
\]
is not presently a proved theorem.  Therefore a formula that inserts $12$ into
a loop coefficient must name the twist assumptions; it cannot be described as
a unique consequence of $\mathbb C\otimes\mathbb H$ alone.

Correspondingly, the repository records two baseline coefficients in different
routes,
\[
 B_0^{\rm loop}=2\pi,
 \qquad
 B_0^{\rm twist}=8\pi.
\]
with the bridge from either baseline to the phenomenologically required
$B\approx46.284$ still open.

\section{Gap G137-B}
The decisive alpha problem can be stated without rhetoric:
\begin{quote}
Starting from the canonical UBT action and its justified field content, derive
the effective reduced coefficient multiplying $n\ln n$, including all
normalisations and multiplicities, and show that the derivation is independent
of the measured value of $\alpha$.
\end{quote}
This is Gap G137-B.  Several historical and research-track routes have produced
interesting modular, eta, Gamma, determinant, and winding structures, but the
current master status records the required insertion into the physical
coefficient as open.

\section{Why it still makes sense to study the theta/Mellin branch}
The failure of the naive chain
\[
 \Theta\to\mathcal M\to\zeta\to\mathcal E\to\mathcal P
\]
where $\mathcal M$, $\mathcal E$, and $\mathcal P$ denote the Mellin transform,
Euler product, and prime sectors, respectively,
to\emph{dynamically} select primes is itself useful.  The Euler product of a
standard theta Mellin transform reflects multiplicative arithmetic already
present in the integer index set.  A stronger mechanism would have to identify
primitive sectors dynamically.

Rational-time theta revivals are one candidate because quadratic Gauss sums
factor over coprime denominators.  In that setting prime powers are genuinely
irreducible building blocks of the revival arithmetic.  Whether the alpha
prime-sector construction can be rewritten in those variables is an explicit
open question, not an established result.

\section{Safe summary}
The statements to remember are:
\begin{enumerate}
 \item the reduced stationary equation and its stability calculus are exact;
 \item the integer/prime structural selection can be studied reproducibly;
 \item $B$ is not yet derived from the UBT action at the required value;
 \item the relevant multiplicity bookkeeping is not fully reconciled;
 \item $\alpha^{-1}=137.036\ldots$ is not a first-principles UBT result today.
\end{enumerate}
This is intentionally less dramatic than older archive text and more useful for
future work, because it tells us exactly which equation must be closed next.
